% ==================== 5 XRD 实验操作总结 ====================
\clearpage
\section{实验操作、注意事项与问题处理}

\subsection{操作要点}

结合本次测试，操作要点按时间顺序归纳如下。

\begin{enumerate}[label=（\arabic*）,leftmargin=2.2em]
  \item \textbf{制样。} 样品在玛瑙研钵中磨细至无颗粒感，填入样品槽后用毛玻璃片压实，上表面与槽沿齐平。槽要填满，只铺薄薄一层会照到槽底；压实要均匀，过松或过紧都会让峰形变差。装好后擦净槽外和台面上的散落粉末。
  \item \textbf{开机。} 先开循环冷却水，水温降到约 19 $^\circ$C 再打开钥匙和 XD-3 软件。管压、管流先设 15 kV、6 mA 预热几分钟，再缓慢升到工作值（txt 文件头记录为 40 kV、30 mA，讲义示例为 36 kV、20 mA）。
  \item \textbf{装样与对中。} 样品台平直地装到测角台上，确认表面没有碰到光路；每次换样都用“校读”校准初始位置，不能沿用上一个样品的读数。
  \item \textbf{测量。} 关好防护门，按 20$^\circ$--80$^\circ$、步宽 0.02$^\circ$ 扫描，过程中不开门、不碰样品台；扫完立即存成 txt，文件名写清样品编号和测试条件。
  \item \textbf{关机。} 管压、管流先降到 15 kV、6 mA，等约 5 min 关 X 射线，约 10 min 后关水，最后左旋钥匙关机。顺序颠倒容易损伤靶和灯丝。
  \item \textbf{记录。} 仪器记录本上写清日期、样品、参数、异常情况和关机时间；原始数据至少备份一份，不要只留在仪器电脑上。
\end{enumerate}

\subsection{注意事项}

\begin{enumerate}[label=（\arabic*）,leftmargin=2.2em]
  \item \textbf{X 射线安全。} 测量前确认防护门已经关严，门联锁和警示灯工作正常；扫描过程中严禁开门。X 射线看不见、闻不到，不能凭感觉判断“有没有出束”。
  \item \textbf{冷却水与靶保护。} 开机前必须确认循环水已经工作、水温降到设定值；关机时先降管压/管流、再关 X 射线、最后关水。冷却水中断或流量不足时，阳极靶会在很短的时间内烧毁。
  \item \textbf{样品状态与择优取向。} 样品要干燥、量足、上表面平整且与槽沿齐平；偏高、偏低或偏离中心都会让峰位整体偏移，低角峰还会变形。样品量太少时 X 射线会照到槽底，背底升高、峰强下降。压样不要用力过猛，片状、针状或层状样品过度压实会让晶粒沿某个方向排列，相对强度偏离 PDF 卡片，必要时改用侧装或轻压。
  \item \textbf{角度范围。} 本次按 20$^\circ$--80$^\circ$ 扫描；起始角不要设得太低、终止角不要设得太高，否则光管和探测器可能相撞。特殊角度范围要经指导教师同意。
  \item \textbf{清洁与记录。} 换样时用酒精棉擦净样品槽和样品台，避免上一样品的残留；实验结束后整理台面，填写仪器使用记录，发现异常及时报告。
  \item \textbf{文件格式。} 仪器导出的 txt 是 GBK 编码，导入 Origin 前要先解决逗号分隔符和表头识别问题；不要直接拖进 Origin 再手工删列，容易把数据行和元数据混在一起。
  \item \textbf{PDF 卡片编号。} 不同数据库版本的编号可能不同（如 71-1167 与 21-1272 都对应锐钛矿），选卡片时要把主要衍射线对一遍，不能只看编号。
\end{enumerate}

\Needspace{6\baselineskip}
\subsection{存在问题与解决办法}

表~\ref{tab:problems} 按“现象—原因—解决办法”列出本次实验和数据处理中遇到的问题。

\begin{longtable}{@{}p{2.4cm}p{3.5cm}p{3.6cm}p{5.7cm}@{}}
  \caption{XRD 实验中常见问题、原因与解决办法}\label{tab:problems}\\
  \toprule
  问题 & 现象 & 原因 & 解决办法 \\
  \midrule
  \endfirsthead
  \multicolumn{4}{@{}l}{\small 续表~\thetable}\\
  \toprule
  问题 & 现象 & 原因 & 解决办法 \\
  \midrule
  \endhead
  \bottomrule
  \endfoot
  原始 txt 导入 Origin 错列 & 数据只剩一列，表头变成 0 kV、0 mA & 逗号被当成分隔符；表头字段数规则误判；编码选错 & 用 \pyA 改成三行表头 + Tab；导入时选 GB2312；按“从新的列开始”导入 \\
  文件头管压/管流记录不全 & 200--800.txt 的 kV/mA 行是 0 kV、0 mA，只有 900.txt 是 40 kV、30 mA & 软件保存时没有读到发生器状态 & 在记录本上记下工作管压/管流；扫描参数界面截图存档；报告中注明来源 \\
  低角背底高、峰弱 & 200 $^\circ$C 在 25$^\circ$ 附近只有宽包；300 $^\circ$C 峰顶噪声明显 & 结晶度低、晶粒极细；样品量少或扫描速度快 & 延长扫描、放慢速度、重新压样；扣背底后谨慎解释，不把宽包当锐峰 \\
  峰位整体偏移 & 所有峰比 PDF 卡片高约 0.1$^\circ$--0.2$^\circ$ & 仪器零点漂移、样品表面高于或低于槽沿 & 每次换样“校读”；压平样品；同批样品尽量一次测完；必要时用标样校正 \\
  两相峰重叠 & 700--800 $^\circ$C 的锐钛矿与金红石峰互相叠加 & 两相共存，部分峰位接近 & 同时叠加两张 PDF 卡片；用多根峰交叉验证；不要只看最强峰 \\
  相对强度偏离卡片 & 某些峰的相对高低与卡片不一致 & 择优取向、颗粒大小、吸收效应 & 充分研磨、轻压、侧装；重复制样；用多根峰的相对强度综合判断 \\
  Jade 检索到错误卡片 & 搜出多个 TiO$_2$ 卡片，峰位对不全 & 锐钛矿、金红石、板钛矿同质多象，卡片版本也不同 & 用化学式和元素过滤；把主要衍射线逐峰核对；本次用锐钛矿 71-1167、金红石 75-1756 \\
  PDF 卡片层位置/比例不对 & 卡片和实验峰横坐标对不上，或卡片太大/太小 & 卡片层的 X 范围没设成 20--80$^\circ$；层大小不合适 & 卡片层 X 轴设为 20--80$^\circ$；主图左 15、上 10、宽 70、高 60，卡片层高 10 左右；双击边框调 Scale 和 Ticks \\
  样品撒漏或污染 & 样品台上有粉末，谱线出现异常峰 & 装样过满、擦拭不彻底 & 用酒精棉擦净样品槽和台面；重新装样；换样时先检查台面 \\
  数据丢失或命名混乱 & 找不到某个样品的数据，或文件名重复 & 未及时保存、命名随意 & 扫描结束立即保存；统一用“温度\_日期\_参数”命名；当天备份 \\
  图件不规范 & 拷屏、坐标不统一、图例遮挡 & 直接用屏幕截图或默认设置 & 用 Origin 重绘；统一坐标范围和单位；导出 600 dpi PNG 或矢量 PDF \\
  X 射线安全与设备 & 门未关严、水未开、升降压过快 & 操作顺序不熟、赶时间 & 按开机/关机清单逐项确认；先水后高压、慢升慢降；异常时立即停束并报告 \\
\end{longtable}

表~\ref{tab:problems} 里有三条需要展开。低角背底高不一定说明仪器有问题：纳米粉体结晶度低、颗粒极细时峰就会又宽又弱，把宽包硬读成锐峰，得到的峰位和强度都不可靠。峰位整体偏移属于系统误差，单看一个样品容易误以为晶格膨胀，把 8 个样品放在一起就能看到所有峰一起平移，这时应先查零点和样品高度，而不是改晶胞参数。两相峰重叠是 700--800 $^\circ$C 数据的主要难点，只看一个最强峰容易判错，把锐钛矿和金红石两张卡片都叠上去、用多根峰交叉确认才稳妥。

软件上的两处麻烦都出在 Origin。一是原始 txt 直接导入时表头错位、数据列合并，用 \pyA 改成三行表头 + Tab 分隔后导入就稳定了。二是卡片层：起初卡片和实验峰的横坐标对不上，把卡片层的 X 轴也设成 20$^\circ$--80$^\circ$、按“主图左 15、上 10、宽 70、高 60，卡片层高 10”的比例调整之后，两套卡片才落到主图下方，图~\ref{fig:pdf-step3}--\ref{fig:pdf-step5} 是这一步的记录。

\subsection{实验心得}

这次实验最直接的感受是 XRD 的数据质量在装样时就基本定了。同样的仪器和扫描参数，样品表面高一点、压得松一点，峰位和峰形就会变；软件里的扣背底、平滑只能修补，补不出原本没有的信息。所以每次换样都要重新校读、重新看样品高度，不能因为上一个样品刚测过就跳过。

另一个体会是 PDF 卡片不能当“自动答案”用。同一个 TiO$_2$ 在 Jade 里能搜出锐钛矿、金红石、板钛矿和多个版本的卡片，只看编号或只看最强峰很容易选错；把主要衍射线逐峰对一遍，再看 700 $^\circ$C 的 27.6$^\circ$ 新峰和 800 $^\circ$C 的两相共存，才能确定哪张卡片对应哪个物相。

数据处理也要留痕。原始 txt、转换后的 \texttt{ready/} 和 \texttt{ready\_norm/}、Origin 工程、导出的图按样品编号和步骤命名，当天备份，否则过几天就分不清哪个文件是原始数据。这次把 txt 转换写成脚本、把最终图件用脚本重绘，也是为了让这一步可以重复。
